% BHU PYQ -- Professional Exam Template
\documentclass[11pt,a4paper]{article}

\usepackage[a4paper,top=2.5cm,bottom=2.5cm,left=2.8cm,right=2.8cm,headheight=14pt]{geometry}
\usepackage{amsmath,amssymb}
\usepackage[T1]{fontenc}
\usepackage[utf8]{inputenc}
\usepackage{lmodern}
\usepackage{microtype}
\usepackage{fancyhdr}
\usepackage{enumitem}
\usepackage{listings}
\usepackage{hyperref}
\usepackage{lastpage}
\usepackage{parskip}

\hypersetup{colorlinks=true,urlcolor=black,linkcolor=black,
  pdftitle={BVCA-103b: Applied Mathematics -- BHU PYQ}}

\pagestyle{fancy}
\fancyhf{}
\renewcommand{\headrulewidth}{0.4pt}
\renewcommand{\footrulewidth}{0.4pt}
\fancyhead[L]{\small   \textit{Banaras Hindu University}}
\fancyhead[R]{\small   \textit{BVCA-103b: Applied Mathematics}}
\fancyfoot[C]{\small Page  \thepage\ of \pageref{LastPage}}

\lstset{
  basicstyle=\small tfamily,
  frame=single,
  breaklines=true,
  columns=flexible,
  keepspaces=true
}


\newlist{parts}{enumerate}{2}
\setlist[parts,1]{label=(\alph*),leftmargin=*,itemsep=4pt,topsep=3pt}
\setlist[parts,2]{label=(\roman*),leftmargin=*,itemsep=2pt,topsep=2pt}

\newcommand{\pts}[1]{\hfill   \textit{[#1]}}


\begin{document}

\begin{center}
  {\large\bfseries BANARAS HINDU UNIVERSITY}\\[2pt]
  {\normalsize B.Voc. Semester I Examination, 2022-23}\\[6pt]
  \rule{0.8\linewidth}{0.5pt}\\[6pt]
  {\large\bfseries Paper: BVCA-103b --- Applied Mathematics}\\[4pt]
  {\normalsize Time: Three Hours \hfill Maximum Marks: 70}
\end{center}

\rule{\linewidth}{0.4pt}

\smallskip



\noindent   \textit{Instructions: Answer all Sections. Candidates are required to solve their answers with accuracy.}

\bigskip

\bigskip


\noindent {\large\bfseries SECTION --- A: Long Answer Questions}\pts{$15  \times 2 = 30$}
\rule{\linewidth}{0.4pt}\smallskip




\noindent   \textit{Instructions: Long answer type questions each have to be solved stepwise.}

\medskip


\noindent   \textbf{Question 1.} Prove that:
\[

\begin{vmatrix}
  1 & 1 & 1 \\
  a & b & c \\
  a^2 & b^2 & c^2
\end{vmatrix} = (a-b)(b-c)(c-a)
\]
\centerline{   \textbf{OR}}

\medskip


\noindent   \textbf{Question 2.} Apply Gauss elimination method to solve the equations:
\[

\begin{aligned}
  x + 4y - z &= -5 \\
  x + y - 6z &= -12 \\
  3x - y - z &= 4
\end{aligned}
\]

\medskip


\noindent   \textbf{Question 3.} Find the eigenvalues and verify the Cayley-Hamilton theorem for the given matrix:
\[
A = 
\begin{pmatrix}
  5 & 4 \\
  1 & 2
\end{pmatrix}
\]
\centerline{   \textbf{OR}}

\medskip


\noindent   \textbf{Question 4.} Define with examples:

\begin{parts}
  \item The following:
  
\begin{parts}
    \item Set
    \item Singleton Set
    \item Finite Set
    \item Infinite Set
    \item Null Set
    \item Subset
  \end{parts}
  \item If $n(A) = 50$, $n(B) = 60$, and $n(C) = 25$, then find $n(A \cup B \cup C)$ if $A, B, C$ are disjoint sets.
\end{parts}

\bigskip


\noindent {\large\bfseries SECTION --- B}
\rule{\linewidth}{0.4pt}\smallskip

Note: Attempt any six out of eight. Short answer type questions to be solved in required steps:\hfill ($5  \times 6 = 30$)

\medskip


\noindent   \textbf{Question 1.} Calculate the number of subsets for the given set $A = \{0, 1, 2\}$ and write all the possible subsets of the given set.

\medskip


\noindent   \textbf{Question 2.} Find the probability of getting Head and Tail both on one toss of a single coin.

\medskip


\noindent   \textbf{Question 3.} Write the characteristic equation of the given matrix:
\[
A = 
\begin{pmatrix}
  1 & 1 & 3 \\
  1 & 5 & 1 \\
  3 & 1 & 1
\end{pmatrix}
\]

\medskip


\noindent   \textbf{Question 4.} Write the statement of Cayley-Hamilton theorem and verify this theorem for an Identity matrix of order $3  \times 3$.

\medskip


\noindent   \textbf{Question 5.} Define Derivation Tree.

\medskip


\noindent   \textbf{Question 6.} Apply Gauss-Seidel iteration method to solve the equations:
\[

\begin{aligned}
  20x + y - 2z &= 17 \\
  3x + 20y - z &= -18 \\
  2x - 3y + 20z &= 25
\end{aligned}
\]
Perform this operation up to two iterations.

\medskip


\noindent   \textbf{Question 7.} Define set of set and write the power set of the given set $A = \{1, 2, 3\}$.

\medskip


\noindent   \textbf{Question 8.} Define the sample space and write the sample space of the outcome of a throw of a dice, and find the probability that we get 6 in a single throw of dice.

\bigskip


\noindent {\large\bfseries SECTION --- C}
\rule{\linewidth}{0.4pt}\smallskip

Note: Attempt all questions. Very short answer type questions or one-word answer type questions:\hfill ($10  \times 1 = 10$)


\begin{enumerate}[label=\arabic*.,leftmargin=*,itemsep=6pt]
  \item Can we calculate the determinant value of a row matrix?
  \item Verify if the Identity matrix of order $2  \times 2$ is a non-singular matrix.
  \item Why is the system of Homogeneous equations called consistent?
  \item Can we refer to a null set as an empty set?
  \item Calculate the value of the given determinant:
  \[
  A = 
\begin{vmatrix}
    1 & 0 & 1 \\
    1 & 0 & 2 \\
    1 & 1 & 2
  \end{vmatrix}
  \]
  \item Write what type of set is $A = \{\}$.
  \item Write the definition of 'Mutually Exclusive Event'.
  \item Calculate the number of terms in the expansion of a $3  \times 3$ determinant.
  \item What is an Upper Triangular Matrix?
  \item Define Function.
\end{enumerate}

\vfill
\rule{\linewidth}{0.4pt}

\begin{center}\small
\href{https://bhu-exam.vercel.app/}{Click here to visit our website}\\[4pt]\href{https://me-aryan.vercel.app/}{Click here to visit our contributor}\\[4pt]\href{https://quashan-me.vercel.app/}{Click here to visit our contributor}
\end{center}
\end{document}